\documentclass[11pt,a4paper]{article}

% =====================================================
% preamble.tex — Configuración común de estilo
% Para enunciados de ejercitación en Tekla Structures
% =====================================================

% --- Codificación e idioma ---
\usepackage[utf8]{inputenc}
\usepackage[T1]{fontenc}
\usepackage[spanish,es-tabla,es-noshorthands]{babel}
\usepackage{lmodern}
\usepackage{microtype}

% --- Geometría de página ---
\usepackage[a4paper,margin=2.5cm,headheight=14pt]{geometry}

% --- Matemática y unidades ---
\usepackage{amsmath,amssymb}
\usepackage{siunitx}
\sisetup{
  output-decimal-marker={,},
  per-mode=symbol,
  group-separator={.},
  group-minimum-digits=4
}

% --- Tablas ---
\usepackage{booktabs}
\usepackage{array}
\usepackage{multirow}
\usepackage{tabularx}
\usepackage{longtable}

% --- Figuras ---
\usepackage{graphicx}
\usepackage{caption}
\usepackage{subcaption}
\usepackage{float}

% --- Gráficos vectoriales ---
\usepackage{tikz}
\usetikzlibrary{arrows.meta,positioning,calc,patterns,shapes.geometric}

% --- Listas ---
\usepackage{enumitem}
\setlist{itemsep=2pt,parsep=2pt}

% --- Color y cajas ---
\usepackage{xcolor}
\definecolor{teklablue}{RGB}{0,87,156}
\definecolor{boxgray}{RGB}{245,245,245}
\definecolor{accentorange}{RGB}{230,120,30}

\usepackage[most]{tcolorbox}

% Caja para comandos / atajos
\newtcolorbox{comando}{
  colback=boxgray,
  colframe=teklablue,
  boxrule=0.5pt,
  arc=2pt,
  left=6pt, right=6pt, top=4pt, bottom=4pt,
  fonttitle=\bfseries,
  title=Comando Tekla
}

% Caja para tips / consejos
\newtcolorbox{tip}{
  colback=yellow!8,
  colframe=accentorange,
  boxrule=0.5pt,
  arc=2pt,
  fonttitle=\bfseries,
  title=Tip
}

% Caja para entregable
\newtcolorbox{entregable}{
  colback=green!5,
  colframe=green!50!black,
  boxrule=0.5pt,
  arc=2pt,
  fonttitle=\bfseries,
  title=Entregable
}

% Caja para advertencia
\newtcolorbox{advertencia}{
  colback=red!5,
  colframe=red!70!black,
  boxrule=0.5pt,
  arc=2pt,
  fonttitle=\bfseries,
  title=Atención
}

% --- Código y comandos en línea ---
\usepackage{listings}
\lstset{
  basicstyle=\ttfamily\small,
  breaklines=true,
  frame=single,
  backgroundcolor=\color{boxgray},
  rulecolor=\color{teklablue}
}

% Comando para resaltar comandos/menús de Tekla en línea
% Uso: \tk{File > Open}  o  \tk{Component Catalog}
\newcommand{\tk}[1]{\textsf{\textbf{\textcolor{teklablue}{#1}}}}

% Comando para nombres de archivo / rutas
\newcommand{\arch}[1]{\texttt{#1}}

% Comando para teclas / shortcuts
% Uso: \key{Ctrl+S}
\newcommand{\key}[1]{\fbox{\footnotesize\texttt{#1}}}

% --- Encabezado y pie de página ---
\usepackage{fancyhdr}
\pagestyle{fancy}
\fancyhf{}
\fancyhead[L]{\small\textit{\tituloEjercicio}}
\fancyhead[R]{\small\thepage}
\fancyfoot[C]{\small Tekla Structures 2022 — Nivel avanzado}
\renewcommand{\headrulewidth}{0.4pt}
\renewcommand{\footrulewidth}{0.4pt}

% --- Hipervínculos y referencias ---
\usepackage{hyperref}
\hypersetup{
  colorlinks=true,
  linkcolor=teklablue,
  urlcolor=teklablue,
  citecolor=teklablue,
  pdfborder={0 0 0}
}
\usepackage[noabbrev,spanish]{cleveref}

% --- Formato de secciones ---
\usepackage{titlesec}
\titleformat{\section}
  {\Large\bfseries\color{teklablue}}
  {\thesection.}{0.5em}{}
\titleformat{\subsection}
  {\large\bfseries\color{teklablue!80!black}}
  {\thesubsection.}{0.5em}{}

% =====================================================
% commands.tex — Variables propias de cada enunciado
% Editar estos campos al inicio para cada ejercicio
% =====================================================

% --- Identificación del ejercicio ---
\newcommand{\tituloEjercicio}{Ejercicio: [TÍTULO]}
\newcommand{\subtituloEjercicio}{[Subtítulo descriptivo]}
\newcommand{\codigoEjercicio}{TEKLA-XXX}
\newcommand{\nivelEjercicio}{Avanzado}
\newcommand{\duracionEjercicio}{[X] horas}

% --- Datos del curso ---
\newcommand{\nombreCurso}{[Nombre del curso]}
\newcommand{\autorEnunciado}{[Docente]}
\newcommand{\fechaEjercicio}{\today}

% --- Versión del software ---
\newcommand{\versionTekla}{Tekla Structures 2022}
\newcommand{\entornoTekla}{[Environment / Configuration]}

% --- Atajos para terminología recurrente ---
\newcommand{\CC}{\tk{Component Catalog}}
\newcommand{\AC}{\tk{Applications \& Components}}
\newcommand{\OB}{\tk{Object Browser}}
\newcommand{\OR}{\tk{Object Representation}}
\newcommand{\OG}{\tk{Object Group}}
\newcommand{\SF}{\tk{Selection Filter}}
\newcommand{\VF}{\tk{View Filter}}


\renewcommand{\tituloEjercicio}{Ejercicio de Filtros sobre Modelo Mixto}
\renewcommand{\subtituloEjercicio}{Selection, View y Object Group filters en proyecto hormigón + acero}
\renewcommand{\codigoEjercicio}{TEKLA-FT-01}

\begin{document}

% =====================================================
% PORTADA
% =====================================================
\begin{titlepage}
\centering
\vspace*{2cm}
{\Huge\bfseries\color{teklablue} \tituloEjercicio \par}
\vspace{0.5cm}
{\Large \subtituloEjercicio \par}
\vspace{2cm}
\rule{0.6\textwidth}{0.4pt}\par
\vspace{1cm}
\begin{tabular}{rl}
\textbf{Código:} & \codigoEjercicio \\
\textbf{Curso:} & \nombreCurso \\
\textbf{Nivel:} & \nivelEjercicio \\
\textbf{Duración estimada:} & \duracionEjercicio \\
\textbf{Software:} & \versionTekla \\
\textbf{Docente:} & \autorEnunciado \\
\textbf{Fecha:} & \fechaEjercicio \\
\end{tabular}
\vfill
\end{titlepage}

\tableofcontents
\newpage

% =====================================================
\section{Objetivos de aprendizaje}
% =====================================================
\begin{itemize}
    \item Construir \SF, \VF\ y \OG\ con criterios simples y combinados (AND / OR).
    \item Aplicar filtros sobre propiedades estándar y \tk{User-Defined Attributes (UDA)}.
    \item Usar \OR\ para asignar colores y transparencias por categoría.
    \item Integrar filtros en planos (\tk{Drawing View Filter}) y reportes (\tk{Template Editor}).
    \item Diseñar un sistema coherente de filtros reutilizable en distintos proyectos.
\end{itemize}

% =====================================================
\section{Modelo de partida}
% =====================================================
Se entrega un modelo mixto (hormigón + acero) con la siguiente composición aproximada:
\begin{itemize}
    \item [Cantidad y tipo de elementos de hormigón]
    \item [Cantidad y tipo de elementos de acero]
    \item [UDAs ya cargados en los objetos]
    \item [Fases (\tk{Phases}) ya definidas]
\end{itemize}

\begin{advertencia}
El modelo de partida \textbf{no debe modificarse} en cuanto a geometría. Todo el ejercicio se resuelve creando filtros y vistas, no nuevos elementos.
\end{advertencia}

% =====================================================
\section{Esquema del modelo}
% =====================================================
\begin{figure}[H]
\centering
\fbox{\parbox[c][6cm][c]{0.8\textwidth}{\centering\textit{[Vista general del modelo entregado]}}}
\caption{Vista isométrica del modelo de partida.}
\end{figure}

% =====================================================
\section{Consigna paso a paso}
% =====================================================

\subsection{Análisis previo del modelo}
\begin{enumerate}
    \item Abrir el modelo entregado y explorar con \OB.
    \item Inventariar clases (\tk{Class}), perfiles, materiales y fases presentes.
    \item Identificar qué UDAs están cargados.
\end{enumerate}

\subsection{Selection Filters}
Crear los siguientes \SF:
\begin{enumerate}
    \item Filtro 1: [criterio simple — ej. todos los elementos de acero]
    \item Filtro 2: [criterio por clase]
    \item Filtro 3: [criterio combinado AND — ej. vigas de acero del nivel 2]
    \item Filtro 4: [criterio combinado OR]
    \item Filtro 5: [criterio con UDA]
    \item Filtro 6: [criterio con expresión regular o wildcard sobre marca]
\end{enumerate}

\begin{comando}
\tk{Select > Selection Filter > New} \\
Categorías relevantes: \tk{Part}, \tk{Assembly}, \tk{Reinforcement}, \tk{Bolt}.
\end{comando}

\subsection{View Filters}
\begin{enumerate}
    \item Crear \VF\ que muestre solo [criterio] y guardar como vista nueva.
    \item Crear \VF\ que oculte temporalmente [criterio].
    \item Combinar \VF\ con \tk{Representation} para distinguir visualmente elementos.
\end{enumerate}

\subsection{Object Representation por filtro}
\begin{enumerate}
    \item Definir colores y transparencias en \OR\ ligados a filtros:
    \begin{itemize}
        \item [Hormigón → gris claro]
        \item [Acero por fase → distintos colores]
        \item [Elementos sin numerar → rojo]
    \end{itemize}
    \item Guardar el esquema con un nombre descriptivo.
\end{enumerate}

\subsection{Object Group filters para reportes}
\begin{enumerate}
    \item Crear un \OG\ específico para reportes de [categoría].
    \item Configurarlo para alimentar un \tk{Report Template} en \tk{Template Editor}.
    \item Generar el reporte y verificar que solo aparezcan los objetos esperados.
\end{enumerate}

\subsection{Filtros en planos}
\begin{enumerate}
    \item Aplicar \tk{Drawing View Filter} en un plano existente para mostrar solo [criterio].
    \item Configurar \tk{Object Level Settings} con filtros para diferenciar la representación de elementos en el plano.
\end{enumerate}

\subsection{Sistema integrado}
Diseñar un conjunto de filtros nombrados con una convención coherente (prefijo, categoría, acción) que cubra:
\begin{itemize}
    \item Filtros de selección por disciplina
    \item Filtros de selección por fase
    \item Filtros de selección por estado de numeración
    \item Filtros de visualización para revisión y para presentación
    \item Filtros para reportes de cómputo
\end{itemize}

\begin{tip}
Una convención de nombres tipo \arch{[disciplina]\_[criterio]\_[acción]} (por ej. \arch{HA\_NIVEL2\_SEL}) facilita la reutilización en proyectos futuros.
\end{tip}

% =====================================================
\section{Caso de uso final}
% =====================================================
Usando exclusivamente los filtros creados, el alumno debe producir, sobre el mismo modelo:
\begin{enumerate}
    \item Una vista 3D solo con la estructura de hormigón.
    \item Una vista 3D solo con la estructura de acero coloreada por fase.
    \item Un reporte de cómputo separado por disciplina.
    \item Un plano filtrado de un nivel específico.
\end{enumerate}

% =====================================================
\section{Entregables}
% =====================================================
\begin{entregable}
\begin{enumerate}
    \item Archivo de modelo \arch{.db1} con todos los filtros creados.
    \item Exportación de los filtros (\arch{.SObjGrp}, \arch{.VObjGrp}, \arch{.objgrp}).
    \item PDF con capturas del caso de uso final (4 vistas + plano).
    \item Reporte de cómputo en formato \arch{.xls} o \arch{.csv}.
    \item Tabla en PDF con la convención de nombres adoptada y descripción de cada filtro.
\end{enumerate}
\end{entregable}

% =====================================================
\section{Criterios de evaluación}
% =====================================================
\begin{table}[H]
\centering
\begin{tabularx}{\textwidth}{Xcc}
\toprule
\textbf{Criterio} & \textbf{Puntaje} & \textbf{Obtenido} \\
\midrule
Selection filters: cobertura y corrección & 20 & \\
Uso correcto de combinaciones AND/OR y UDAs & 15 & \\
View filters y object representation & 15 & \\
Object Group e integración con reportes & 15 & \\
Aplicación de filtros en planos & 10 & \\
Coherencia y reutilización del sistema de nombres & 15 & \\
Completitud del caso de uso final & 10 & \\
\midrule
\textbf{Total} & \textbf{100} & \\
\bottomrule
\end{tabularx}
\end{table}

% =====================================================
\section{Tips y errores comunes}
% =====================================================
\begin{tip}
[Tip sobre uso de wildcards o expresiones regulares en filtros]
\end{tip}

\begin{advertencia}
[Error frecuente: por ejemplo, confundir Selection Filter con View Filter]
\end{advertencia}

% =====================================================
\section{Referencias y documentación}
% =====================================================
\begin{itemize}
    \item Tekla User Assistance — Filters: \url{https://support.tekla.com/}
    \item [Otras referencias]
\end{itemize}

\end{document}
